\subsection{Imputation: Problem Formulation and Characteristics}
\label{sec:imputation}

Financial time-series imputation is fundamentally formulated as a bidirectional, closed-format generation task. Given a multivariate sequence $\mathbf{X} \in \mathbb{R}^{T \times d}$ and a corresponding binary indicator mask $\mathbf{M} \in \{0, 1\}^{T \times d}$ where $m_{t,i} = 1$ denotes a missing value, the objective is to learn a conditional distribution $p(\mathbf{X}_{miss} | \mathbf{X}_{obs}; \theta)$ to accurately reconstruct the missing elements $\mathbf{X}_{miss} = \mathbf{X} \odot \mathbf{M}$ conditioned strictly on the observed elements $\mathbf{X}_{obs} = \mathbf{X} \odot (\mathbf{1} - \mathbf{M})$.

In financial contexts, the missingness mechanisms—missing completely at random (MCAR), missing at random (MAR), and missing not at random (MNAR), formalized by Rubin~\citep{rubin1976inference}—present unique challenges. MNAR is particularly prevalent due to structural market events such as sudden trading halts during market crashes or asynchronous trading hours across global exchanges, a setting related to non-synchronous observation of multivariate financial processes~\citep{hayashi2005covariance}. Consequently, imputation models must conquer two primary domain-specific hurdles. First, models must manage \textit{contiguous missingness}, where extended data outages corrupt long temporal windows. Second, models must maintain strict \textit{cross-variable consistency} without introducing look-ahead bias, which is critical to ensure that historical backtests of trading strategies remain structurally valid.

\subsection{Imputation: Statistical, Interpolation, and Machine-Learning Baselines}

Early methodologies treated financial imputation primarily as a matrix-completion~\citep{candes2009exact} or interpolation problem rooted in explicit statistical assumptions. The classic Chow-Lin method \citep{chow1971} remains a foundational statistical approach for temporal disaggregation, frequently integrated with LASSO regression~\citep{tibshirani1996regression} to estimate high-frequency (e.g., monthly) values from low-frequency (e.g., quarterly) observations using auxiliary market indicators. Later, machine-learning techniques such as $k$-nearest-neighbour imputation~\citep{troyanskaya2001missing}, multiple imputation by chained equations (MICE)~\citep{vanbuuren1999multiple}, and Random Forests~\citep{breiman2001random}, including missForest~\citep{stekhoven2012}, offered flexible, non-parametric alternatives for tabular data imputation. An adversarial learning network using heterogeneous data sources for FinTech applications~\citep{NDDP} is included here as a related data-driven baseline. Operating in a continuous functional space, kernel ridge regression~\citep{saunders1998dual,aronszajn1950theory} has been carried into fixed-income curve construction~\citep{smith2020}, where noisy bond prices are projected into a reproducing kernel Hilbert space (RKHS) to recover the discount curve. In an extensive Swiss government-bond study, the kernel-ridge estimator proved robust, transparent, and reproducible, outperforming the regulatory Swiss Solvency Test and central-bank benchmarks both in- and out-of-sample; the authors also document the intrinsic difficulty of extrapolating the yield curve beyond observed maturities, a limitation shared by every estimator they compared.

As datasets expanded, traditional machine learning models such as $k$-nearest neighbours~\citep{troyanskaya2001missing}, MICE~\citep{vanbuuren1999multiple}, and Random Forest/missForest imputation~\citep{breiman2001random,stekhoven2012} were widely adopted. Random Forest imputation, in particular, leverages ensemble decision trees to dynamically predict missing values by capturing complex, non-linear feature interactions. Evidence from corporate-distress prediction indicates that tree-based ensembles can remain competitive with, and in some sparse-data settings outperform, deep learners~\citep{sideras2024bankruptcy}; the precise magnitude of the advantage is dataset-dependent.

\subsection{Imputation: Tensor Factorization and Deep Sequential Approaches}

To accommodate higher-dimensional non-linearities and multivariate temporal dependencies, research progressed toward deep interpolation and tensor decomposition networks. Deep sequential fusion architectures, such as bidirectional GRU--LSTM frameworks, capture bidirectional temporal flows using gated recurrent units~\citep{cho2014learning} within a bidirectional recurrent architecture~\citep{schuster1997bidirectional}, and employ a group method of data handling (GMDH) polynomial decoder~\citep{ivakhnenko1971polynomial} to reconstruct complex missing patterns in multivariate series.

Concurrently, tensor factorization models emerged as highly efficient solutions for structural missingness. NMTucker~\citep{varolgunes2023nmtucker}, building on Tucker decomposition~\citep{tucker1966some}, replaces the multi-linear core with neural network-based reconstructions arranged in a Matryoshka (nested) hierarchy, so that factorizations at successively larger tensor ranks regularize one another and remain stable on the highly sparse time$\times$firm$\times$attribute tensors typical of firm fundamentals. Expanding upon this, RegTensor~\citep{zhou2023regtensor} targets the same coupled tensor-completion problem with an explicit anti-overfitting design: an embedding-learning module for each tensor order, an MLP that models non-linear interactions among the resulting embeddings, and an orthogonal-regularization module that penalizes correlated factors as the tensor rank grows. The regularizer takes the form
\begin{equation}
\mathcal{L}_{RegTensor} = \mathcal{L}_{recon} + \lambda \|\mathbf{U}^\top \mathbf{U} - \mathbf{I}\|_F^2
\end{equation}
where $\mathbf{U}$ represents the temporal factor matrix and $\lambda$ controls the strength of the orthogonality penalty. While computationally efficient, these models occasionally fail to capture the high-frequency volatility jumps inherent to tick-level financial data.

\subsection{Imputation: Adversarial and Generative Imputation Dynamics}

To overcome the smoothing effects typical of standard interpolators, Generative adversarial networks (GANs)~\citep{goodfellow2014gan,yoon2018gain} formulate missing data recovery as a competitive game. However, standard GANs frequently suffer from severe distributional incoherence—commonly referred to as "drift"—during contiguous data outages. To directly address this, AttnWGAIN~\citep{wang2025attn} replaces the convolutional generator with a stack of diagonally masked attention blocks, forcing the generator to capture temporal and cross-feature correlations jointly while the discriminator scores observed against imputed elements. Introduced for IoT sensor streams, the design transfers naturally to financial panels that exhibit the same contiguous-outage structure. Similarly, ImputeGAN~\citep{qin2023imputegan} couples a GAN with an iterative gradient-based refinement strategy, improving both the stability and the generalisability of the completed values without assuming the missingness pattern in advance.

To target the harder case of long contiguous gaps, Wang et al.~\citep{wang2024cwgain} train a conditional Wasserstein GAN with a gradient penalty~\citep{arjovsky2017wgan,gulrajani2017wgangp} explicitly for continuous missing-value blocks. By conditioning the adversarial game on the observation mask $\mathbf{M}$ and enforcing Lipschitz continuity on the critic, the resulting model bounds generative drift so that the imputed blocks remain consistent with the heavy-tailed distribution of the surrounding observable series.

\subsection{Imputation: Diffusion and Foundation Models for Missing Data}

Score-based diffusion models~\citep{ho2020ddpm,song2019generative,song2021score} have become a leading approach for probabilistic imputation by treating missing value reconstruction as a conditional generative process. These architectures define a forward process that adds Gaussian noise strictly to the missing regions (or the entire sequence, masked subsequently), and optimize a reverse denoising step conditioned on $\mathbf{X}_{obs}$:
\begin{equation}
\mathcal{L}_{diff} = \mathbb{E}_{\mathbf{x}_0, \boldsymbol{\epsilon}, t}[\|\boldsymbol{\epsilon} - \boldsymbol{\epsilon}_\theta(\mathbf{x}_t, t, \mathbf{X}_{obs})\|^2]
\end{equation}
A central architecture in this domain is CSDI \citep{tashiro2021csdi}, which utilizes a dual 2D-attention mechanism to seamlessly integrate temporal dynamics and cross-feature correlations. Subsequent works such as LSSDM, Score-CDM, and SaSDim \citep{LSSDM, Score-CDM, SaSDim} further advanced this paradigm. To eliminate the "boundary disharmony" problem—unrealistic volatility jumps at the borders of observed and imputed regions—advanced diffusion models incorporate explicit consistency constraints across temporal and spatial horizons, as well as across cross-variable boundaries.

Furthermore, models like LSCD \citep{fons2025lscd} and SPDM \citep{SPDM} employ frequency-domain transformations to map 1D time series into 2D periodic representations. This isolates stable macroeconomic cyclical components from high-frequency market microstructure noise during the denoising phase. Expanding this to transformer scales, TimeDiT \citep{cao2024timedit} introduces a diffusion Transformer (DiT)~\citep{peebles2023scalable} utilizing a unified masking strategy, uniquely enabling fine-tuning-free model editing. This allows quantitative analysts to directly inject external physical constraints or market heuristics into the sampling trajectory.

Simultaneously, large language models (LLMs), building on few-shot language-model capabilities~\citep{brown2020language}, have been adapted for heterogeneous tabular and sequential data imputation. LLM-Forest \citep{chen2023llmforest} pioneers an ensemble approach utilizing graph-augmented prompts. It constructs bipartite graphs to map data entries to feature values, utilizing hierarchical merging to retrieve nearest neighbors. These neighbors serve as few-shot prompts for a forest of LLMs, resolving the final imputed value via confidence-weighted voting. Additionally, frameworks like MTabGen \citep{liu2023mtabgen} bridge diffusion dynamics and mixed-type tabular formats by running parallel forward processes for categorical and continuous numerical features, unifying them via columnar embeddings.

\subsection{Imputation: Emerging Trends}

Recent literature indicates a substantial shift toward handling complex, non-random missing patterns through specialized extensions of these generative frameworks. As documented in recent methodological surveys~\citep{wang2024mtsisurvey,richter2024survey}, the field is rapidly moving away from isolated point-interpolators toward unified, context-aware sequence modelers. We highlight three critical emerging trends driving this evolution.

The first trend concerns unified multi-task architectures: instead of treating imputation as an isolated preprocessing step, modern models increasingly solve related financial tasks simultaneously to improve representational learning. For example, coupled diffusion models have been developed that jointly perform forecasting and imputation, leveraging shared latent representations to enforce temporal consistency. To enhance structural robustness, dual-attention networks simultaneously capture local and global dependencies, while temporal graph networks (TGNs)~\citep{rossi2020temporal} map dynamic cross-asset spillovers. Regardless of architecture, adversarial learning has become a dominant recipe for these unified pipelines~\citep{richter2024survey}. Furthermore, recent multi-modal frameworks extend this paradigm beyond numerical data, integrating textual sentiment to inform the imputation of missing price data~\citep{T2S}.

The second trend focuses on high-frequency and microstructure imputation: it centers on the granular complexities of high-frequency trading (HFT) and Limit Order Book (LOB) data. Standard daily-level imputers often fail at the microstructural level; thus, specialized deep learning architectures and benchmark suites have been proposed that strictly respect LOB queue dynamics and volume imbalances~\citep{lobench2025}. To bridge the gap between generative fidelity and the stringent latency requirements of HFT systems, scalable probabilistic imputation techniques are being developed that bypass the computational overhead of iterative diffusion without sacrificing distributional shape~\citep{richter2024survey}.

The third trend concerns regime-awareness and stress-testing protocols: the evaluation of imputation models has matured to account for market non-stationarity. Recognizing that structural breaks often coincide with missing data (e.g., trading halts), dynamic, regime-aware score-based models have been introduced specifically to adapt to shifting market volatility~\citep{wang2024mtsisurvey}. Crucially, a growing body of work advocates for \textit{stratified stress-testing evaluation protocols}. Imputation fidelity must be validated during simulated market crashes rather than strictly during stationary periods, ensuring that imputed data maintains structural integrity when risk managers need it most. Complementing this evaluation, the pervasive issue of look-ahead bias in deep learning imputers has been investigated, motivating strict cross-validation guidelines to ensure that reconstructed historical sequences do not inadvertently leak future market information into backtesting pipelines~\citep{wang2024mtsisurvey}. Such release-timing effects are empirically detectable in macro-driven markets~\citep{frino2018macroeconomic}.

\subsection{Imputation: Research Synthesis}

The landscape of financial missing data recovery is defined by a rigorous trade-off between computational efficiency and generative fidelity, as synthesized in Table~\ref{tab:imputation_comparison}. The empirical case for diffusion-based imputation is, however, less settled than the tabular taxonomy above might suggest. The most-cited diffusion imputer, CSDI~\citep{tashiro2021csdi}, reports a 40--65\% improvement over probabilistic baselines and a 5--20\% improvement over deterministic baselines---but on \emph{healthcare and environmental} datasets, not financial. Direct head-to-head benchmarks that put diffusion imputers (CSDI, TimeDiT, LSCD) on the same financial dataset as GAN imputers (GAIN, AttnWGAIN, ImputeGAN, CWGAIN-GP) are not standard in the literature, and the few within-paper comparisons that exist (e.g.\ Score-CDM vs GAIN on air-quality or electricity data) do not transfer cleanly to the financial regime of contiguous missingness and regime-switching volatility. The honest reading is therefore: diffusion models are the most promising current direction for high-fidelity imputation under contiguous missingness, but the claim that they ``stand unrivaled'' for financial imputation specifically rests on architectural plausibility and on transferred evidence rather than on a published financial benchmark. Tensor factorization and tree-based methods remain preferable when the constraint is large-scale historical-database throughput rather than peak reconstruction quality.
